\documentclass{beamer}

\def\Tiny{\fontsize{6pt}{6pt}\selectfont}
\def\supertiny{\fontsize{4pt}{4pt}\selectfont}

\mode<presentation>
{
  \usetheme{Warsaw}
  % \setbeamercovered{transparent}
  \usecolortheme{crane}
}

\usepackage{graphicx, ifthen, listings, fancyvrb}

\usepackage[czech]{babel}
% \usefonttheme{professionalfonts}
\usepackage{times}
\usepackage{amsmath}
\usepackage[utf8]{inputenc}
\usepackage{wrapfig}

\usepackage[T1]{fontenc}

\lstset{ basicstyle=\tiny, stringstyle=\ttfamily, showstringspaces=false }

\everymath{\displaystyle}

\setbeamerfont{frametitle}{size=\large}
\setbeamerfont{subsection in toc}{size=\scriptsize}

\makeatletter\newenvironment{blackbox}{%
   \begin{lrbox}{\@tempboxa}\begin{minipage}{0.95\columnwidth}}{\end{minipage}\end{lrbox}%
   \colorbox{black}{\usebox{\@tempboxa}}
}\makeatother

\title[IMF (11)]{Informatika pro moderní fyziky (9)\\ tvorba vektorové grafiky, formát JSON, použití cizích API}

\author[Franti\v{s}ek HAVL\r{U}J, ORF ÚJV Řež]{Franti\v{s}ek HAVL\r{U}J\\{\scriptsize \emph{e-mail: haf@ujv.cz}}}

\date{akademický rok 2016/2017\\30. listopadu 2016}

\institute[ORF ÚJV Řež]
{ÚJV Řež\\oddělení Reaktorové fyziky a podpory palivového cyklu}

\AtBeginSection[]
{
\begin{frame}<beamer>
\frametitle{Obsah}
\tableofcontents[currentsection,hideothersubsections]
\end{frame}
}

\begin{document}

\begin{frame}
  \titlepage
\end{frame}

\begin{frame}
  \tableofcontents
\end{frame}

\section{Tvorba obrázků}

\begin{frame}{Zadání dnešní úlohy}
  \begin{itemize}
    \item pro zadanou textovou mapu AZ VR1 potřebuju udělat hezký obrázek
    \item co druh, to barvička, rozumně zacházet s odstíny (palivo různě modré, R/B/E tyče různě červené, zelené, fialové)
  \end{itemize}
\end{frame}

\begin{frame}{Jak na obrázky}
  \begin{itemize}
    \item pěkný formát na tvorbu vektorových obrázku je SVG (Scalable Vector Graphics)
    \item je to dobrá věc především na internet -- všechny prohlížeče ho umí
    \item stejně jako HTML je postaven na XML
  \end{itemize}
\end{frame}

\begin{frame}[fragile]{Jednoduchý příklad}
  \tiny
  \begin{verbatim}
    <svg width="320" height="320" xmlns="http://www.w3.org/2000/svg" version="1.1">
      <rect x="0.0" y="0.0" width="40.0" height="40.0" fill="blue" />
      <rect x="40.0" y="0.0" width="40.0" height="40.0" fill="red" />
      <rect x="0.0" y="40.0" width="40.0" height="40.0" fill="green" />
      <rect x="40.0" y="40.0" width="40.0" height="40.0" fill="yellow" />
    </svg>
  \end{verbatim}
  \includegraphics[width=0.15\textwidth]{example}
\end{frame}

\begin{frame}{SVG -- co a jak}
  \begin{itemize}
    \item souřadný systém z levého horního rohu
    \item je potřeba udat celkovou šířku a výšku
    \item zatím nám stačí obdélník -- tag \texttt{rect}
    \item pozor, je to striktní XML, tedy je nutné \texttt{rect} tag uzavřít (!)
    \item vyzkoušejte -- nejdřív jen tak, potom vygenerovat 8x8 mapu (zatím klidně prázdnou)
  \end{itemize}
\end{frame}

\begin{frame}{Další SVG chytrosti}
  \begin{itemize}
    \item kromě \texttt{rect} se bude hodit také \texttt{text}
    \item jako text se zobrazí obsah příslušného elementu
    \item opět použiju atributy \texttt{x}, \texttt{y} (levý dolní roh) a můžu přihodit \texttt{text-anchor=``middle''}, aby to byl dolní prostředek
  \end{itemize}
\end{frame}

\begin{frame}{Postup}
  \begin{itemize}
    \item načtu ze souboru třeba do 2D pole
    \item budu mít hash s barvičkama
    \item vykreslím do SVG
  \end{itemize}
\end{frame}

\begin{frame}[fragile]{Vylepšení: konfigurovatelný}
  \begin{itemize}
    \item načítejme barvy ze souboru!
    \item s hvězdičkou: RRGGBB, \#RRGGBB i R,G,B
  \end{itemize}
\begin{verbatim}
F4  ff0000
R1  00ffff
R2  00eeee
\end{verbatim}
\end{frame}

\section{Persistence dat a formát JSON}

\begin{frame}{Ukládání strukturovaných dat}
  \begin{itemize}
    \item často mám data ve formě struktury (kombinace hash+pole, různý stupeň vnoření)
    \item z různých důvodů můžu chtít data uložit na disk a pak je znovu načítat
    \item (zejména efektivita zpracování, případně data z externích/webových zdrojů)
    \item bylo by dobré mít možnost uložit a načíst rovnou celý hash
    \item odpověď jsou strukturované metaformáty - YAML, XML, JSON
  \end{itemize}
\end{frame}

\begin{frame}{Práce s JSON}
  \begin{itemize}
    \item v Ruby je k mání knihovna -- \texttt{require `json'}
    \item generování JSON: \texttt{hash.to\_json}
    \item čtení JSON: \texttt{JSON[data]}
  \end{itemize}
\end{frame}

\begin{frame}{Příklad -- výsledky běhu kolem rybníka}
  \begin{itemize}
    \item dva soubory -- \texttt{ages.csv}, \texttt{times.csv}
    \item chci v jednom skriptu (tasku) načíst, spárovat a uložit
    \item a v jiném už rovnou načíst zpracovaná data
    \item a vypsat tabulku výsledků včetně ročníků narození
    \item pozor! JSON nezná symboly, uloží se jako řetězce
  \end{itemize}
\end{frame}

\section{Použití cizích API}

\begin{frame}{K čemu to?}
  \begin{itemize}
    \item spousta informací na webu je poskytována ve strojově čitelné formě
    \item API -- rozhraní mezi aplikacemi
    \item s využitím webových služeb naše možnosti exponenciálně rostou (počasí, doprava, mapy, atd atd.)
    \item spousta věcí se dá udělat jako \emph{mashup} -- sice nic neumím, ale umím to dát dohromady
  \end{itemize}
\end{frame}


\begin{frame}{Typy / formáty}
  \begin{itemize}
    \item URL -- rovnou dostanu např. obrázek po zadání správného URL
    \item XML -- velmi obecný, ale komplikovaný formát (``vypadá jako HTML'')
    \item JSON -- velmi jednoduchý a kompaktní formát, vyvinutý pro JS (v podstatě jen číslo, řetězec, pole, hash)
  \end{itemize}
\end{frame}


\begin{frame}{URL API -- google maps}
  \begin{itemize}
    \item stačí správně vymyslet
    \item pozor na usage limits (v produkci je nutné lokální cache...)
    \item QR platba: \\
    {\tiny \url{http://qr-platba.cz/pro-vyvojare/restful-api/\#generator-czech-image}}
    \item Google Maps static API: \\
    {\tiny \url{https://developers.google.com/maps/documentation/staticmaps/}}
  \end{itemize}
\end{frame}


\begin{frame}{Jednoduchý mashup: mapa o-závodů}
  \begin{itemize}
    \item ORIS API -- http://oris.orientacnisporty.cz/API
    \item úkol: vypišme kalendář MTBO závodů v roce 2016
    \item {\tiny \url{http://oris.orientacnisporty.cz/API/?format=json\&method=getEventList\&sport=3\&datefrom=2016-01-01\&dateto=2016-12-31}}
    \item {\tiny \url{https://developers.google.com/maps/documentation/javascript/examples/map-simple}}
    \item API klíč -- v praxi si musíte pořídit vlastní (ale nic to nestojí)
    \item AIzaSyC6\_NpMTN-8olOHWzaiwjeZ8eu\_J3XI1lM
  \end{itemize}
\end{frame}

\begin{frame}{A to je vše, přátelé!}
  \begin{center}
    \includegraphics[width=0.8\textwidth]{looney_tunes}
  \end{center}
\end{frame}

\end{document}
